\documentclass[12pt, twoside]{article}
\input{../preamble}

\fancyhead[LO]{La Scuola d'Italia / Dr.\ Huson / IB Math (b) \\* 17 April 2026}
\fancyhead[RO]{\fbox{\begin{minipage}{6cm}\footnotesize First \& Last name:\\[0.5cm]\end{minipage}}}
\fancyfoot[C]{\thepage}

\begin{document}
\subsubsection*{6.8 Quiz: Polynomials Challenge (no calculator)}

\begin{enumerate}[itemsep=0.15cm]

%--- Problem 1: Factored form quadratic [7 marks] ---
%--- Modeled after AA SL 2021 Nov P1 Q01 ---
\item[1.] [Maximum mark: 7]

Consider the function $f(x) = -2(x + 3)(x - 5)$, for $x \in \mathbb{R}$.

\begin{enumerate}[itemsep=0.15cm]
  \item For the graph of $f$
  \begin{enumerate}[itemsep=0.1cm]
    \item[(i)] write down the $x$-intercepts; \hfill [2]
    \item[(ii)] find the coordinates of the vertex. \hfill [3]
  \end{enumerate}
\end{enumerate}

The function $f$ can be written in the form $f(x) = -2(x - h)^2 + k$.

\begin{enumerate}[itemsep=0.15cm]
  \item[(b)] Write down the value of $h$ and the value of $k$. \hfill [2]
\end{enumerate}

\begin{tikzpicture}
  \draw (-0.6,0) rectangle (11,13);
  \foreach \y in {1,2,...,12} \draw[dotted] (0,\y)--(10.5,\y);
\end{tikzpicture}

\newpage

%--- Problem 2: Tangent line via discriminant [7 marks] ---
\item[2.] [Maximum mark: 7]

The line $y = mx$ passes through the origin and is tangent
to the parabola $y = x^2 + 16$.

\begin{enumerate}[itemsep=0.15cm]
  \item Show that at any point of intersection,
  $x^2 - mx + 16 = 0$. \hfill [1]
  \item Find the discriminant of $x^2 - mx + 16 = 0$
  in terms of $m$. \hfill [1]
  \item Find the values of $m$ for which the line is tangent
  to the parabola. \hfill [2]
  \item For each value of $m$, find the coordinates of
  the point of tangency. \hfill [3]
\end{enumerate}

\begin{tikzpicture}
  \draw (-0.6,0) rectangle (11,13);
  \foreach \y in {1,2,...,12} \draw[dotted] (0,\y)--(10.5,\y);
\end{tikzpicture}

\newpage

%--- Problem 3: Completing the square with a ≠ 1 [7 marks] ---
\item[3.] [Maximum mark: 7]

Consider the function $f(x) = 3x^2 + 6x - 2$.

\begin{enumerate}[itemsep=0.15cm]
  \item Write $f(x)$ in the form $a(x - h)^2 + k$. \hfill [3]
  \item Write down the coordinates of the vertex. \hfill [1]
  \item Find the discriminant of $f(x) = 0$. \hfill [1]
  \item State the number of $x$-intercepts of the graph
  and justify your answer. \hfill [2]
\end{enumerate}

\begin{tikzpicture}
  \draw (-0.6,0) rectangle (11,14);
  \foreach \y in {1,2,...,13} \draw[dotted] (0,\y)--(10.5,\y);
\end{tikzpicture}

\newpage

%--- Problem 4: Cubic polynomial and factor theorem [8 marks] ---
\item[4.] [Maximum mark: 8]

Let $p(x) = 2x^3 - 3x^2 - 11x + 6$.

\begin{enumerate}[itemsep=0.15cm]
  \item Show that $(x - 3)$ is a factor of $p(x)$. \hfill [2]
  \item Hence express $p(x)$ in the form
  $(x - 3)(ax^2 + bx + c)$. \hfill [2]
  \item Factorise $p(x)$ completely into three
  linear factors. \hfill [2]
  \item Write down all solutions to $p(x) = 0$. \hfill [1]
  \item Write down the $y$-intercept of the graph
  $y = p(x)$. \hfill [1]
\end{enumerate}

\begin{tikzpicture}
  \draw (-0.6,0) rectangle (11,13);
  \foreach \y in {1,2,...,12} \draw[dotted] (0,\y)--(10.5,\y);
\end{tikzpicture}

\end{enumerate}
\end{document}
